\documentclass[french]{book}
\usepackage{teach}
\usepackage{verbatim}
\usepackage{amsmath,amsfonts,amssymb}
\usepackage{pgfplots}
\usepackage{tikz}
\usepackage{mwe}
\RequirePackage{graphicx}
\RequirePackage{ifpdf}
 \ifpdf
   \DeclareGraphicsRule{*}{mps}{*}{}
 \fi




\RequirePackage{fancyhdr}
\RequirePackage{titlesec}

  \usepackage{fancyhdr}

  \usepackage{amsmath}
  \usepackage{amssymb}
  \usepackage{amsfonts}
  \usepackage{amsthm}
  \usepackage{mathrsfs}

  \usepackage{array}
  \usepackage{multicol}
  \setlength{\multicolsep}{3pt}

  \usepackage{graphicx}
  \graphicspath{{Figures/}}
  \usepackage{pstricks,pst-plot,pst-text,pst-tree,pst-eucl}
  \usepackage[all]{xy}

  \usepackage{fancybox}
  \usepackage{color}
    \usepackage{stmaryrd}
\usepackage{etex}
%\usepackage{preambuleTrm}
%\usepackage{enumitem}
%\usepackage{tikz}
%
\usetikzlibrary{calc}
%\usepackage[xcas]{pro-courbes}
%\usepackage[autolanguage]{numprint}
\usepackage[french,vlined]{algorithm2e}
\usepackage{cclicenses}
%\usepackage{cclicence}
%\usepackage{docacompleter}

\newcommand{\cacheb}[1]{#1}
\newcommand{\cacheg}[1]{#1} 
\newcommand{\cacher}[1]{#1}
\newcommand{\cachegg}[1]{#1} 

% \newcommand{\cacheb}[1]{\dotfill\textcolor{white}{#1}\textcolor{white}{#1}}
% \newcommand{\cacheg}[1]{\dotfill\textcolor{0.9white}{#1}\textcolor{0.9white}{#1}} 

% \newcommand{\cachegg}[1]{\dots\dots\textcolor{0.8white}{#1}\textcolor{0.8white}{#1}} 
\newcolumntype{A}{>{\arraybackslash}X}

\usepackage{textcomp}
\usepackage{mathtools}
%\usepackage[xcas]{pro-statdiv}
%  \usepackage[frenchb]{babel}
  \usepackage{hyperref}

\usepackage{subfig}
\pgfplotsset{compat=newest}
\usepackage[all]{xy}
%\usepackage{geometry}
%\geometry{top=1cm, bottom=1cm, left=2cm, right=2cm}
\usepackage[utf8]{inputenc}
\usepackage[french]{babel}
\redefinecolor{thm}{title}{bgcolor}{alizarin}
\redefinecolor{prop}{title}{bgcolor}{meatbrown}
\redefinecolor{defn}{title}{bgcolor}{olivine}
\definecolor{alizarin}{rgb}{0.82, 0.1, 0.26}
\definecolor{meatbrown}{rgb}{0.9, 0.72, 0.23}
\definecolor{olivine}{rgb}{0.6, 0.73, 0.45}
\usepackage{mathrsfs}
%%
\newcommand{\R}{\mathbb {R}}
%\newcommand{\C}{\mathds {C}}
\newcommand{\N}{\mathbb {N}}
\newcommand{\Z}{\mathbb {Z}}
\newcommand{\Q}{\mathbb {Q}}
\newcommand{\D}{\mathbb {D}}
%%
\newcommand{\se}{\geq}
\newcommand{\ie}{\leq}
\newcommand{\tm}{\times}
\newcommand\binomialCoefficient[2]{%
    % Store values 
    \c@pgf@counta=#1% n
    \c@pgf@countb=#2% k
    %
    % Take advantage of symmetry if k > n - k
    \c@pgf@countc=\c@pgf@counta%
    \advance\c@pgf@countc by-\c@pgf@countb%
    \ifnum\c@pgf@countb>\c@pgf@countc%
        \c@pgf@countb=\c@pgf@countc%
    \fi%
    %
    % Recursively compute the coefficients
    \c@pgf@countc=1% will hold the result
    \c@pgf@countd=0% counter
    \pgfmathloop% c -> c*(n-i)/(i+1) for i=0,...,k-1
        \ifnum\c@pgf@countd<\c@pgf@countb%
        \multiply\c@pgf@countc by\c@pgf@counta%
        \advance\c@pgf@counta by-1%
        \advance\c@pgf@countd by1%
        \divide\c@pgf@countc by\c@pgf@countd%
    \repeatpgfmathloop%
    \the\c@pgf@countc%
}
\makeatother
%\newcommand{\ell}{l}
%\newcommand{\ell'}{l'}
%%
\newcommand{\limn}{\lim_{n\rightarrow +\infty}}
\newcommand{\limo}[1][x]{\ds\lim_{#1\rightarrow 0}}
\newcommand{\limite}[2][x]{\ds\lim_{#1\rightarrow #2}}
\newcommand{\limpinf}[1][x]{\ds\lim_{#1\rightarrow +\infty}}
\newcommand{\limminf}[1][x]{\ds\lim_{#1\rightarrow -\infty}}
\newcommand{\limiteg}[2][x]{\ds\lim_{#1\xrightarrow{<}#2}}
\newcommand{\limited}[2][x]{\ds\lim_{#1\xrightarrow{>}#2}}
\newcommand{\co}[2]{C_#2^#1}
%%
\newcommand{\vect}[1]{\overrightarrow{#1}}
\newcommand{\oij}{(\textit{O},{\overrightarrow{i}},{\overrightarrow{j}})}
%
\newcommand{\cp}[2]{%
        \begin{pmatrix}
        #1\\
        #2
        \end{pmatrix}%
        }
%
\newcommand{\calb}{\mathscr{B}}
\newcommand{\calc}{\mathscr{C}}
\newcommand{\cald}{\mathscr{D}}
\newcommand{\cale}{\mathscr{E}}
\newcommand{\calf}{\mathscr{F}}
\newcommand{\calp}{\mathscr{P}}
\newcommand{\cals}{\mathscr{S}}
\newcommand{\calt}{\mathscr{T}}
%
\newcommand{\suite}[1][u]{\left( #1_{n} \right)_{n\in\N}}
\newcommand{\suitep}[1][u]{\left( #1_{n} \right)_{n\in\N^{*}}}
\usetikzlibrary{arrows}


\begin{document}
\tableofcontents


\chapter{Ensembles, cardinaux et dénombrements}
\section{Ensembles, $k$-uplet et cardinaux}
Soit $A$ un ensemble et soit $x,y,z\in A$ trois \textbf{éléments} de $A$. 

\begin{defn}
On dit que $(x;y)$ est un \textbf{couple} de l'ensemble $A$ et que $(x;y;z)$ est un \textbf{triplet} de l'ensemble $A$.\\

\par De manière générale, soit $a_1,a_2,\cdots,a_n \in A$, $n$ éléments de $A$, on dit que  $(a_1;a_2;\cdots;a_n) $ est un \textbf{$n$-uplet} d'éléments de  l'ensemble $A$.
\end{defn}

\subsection{Produit cartésien d'ensembles}
Soit $A,B$ deux ensembles.

\begin{defn}
On note $A \times B$ le\textbf{ produit cartésien de $A$ et $B$}. Il contient tous les éléments du type $(a;b)$ avec $a\in A$ et $b\in B$. Autrement dit,
$$A\times B=\{(a;b) / a \in A, b \in B \} $$ 
\end{defn} 

\begin{rem}
Attention, contrairement à la multiplication, le produit cartésien n'est pas commutatif, ils ont même symbole mais pas même signification. 

$$A \times B \neq B \times A$$ 
\end{rem}

\begin{defn}
On note $A^k$ l'ensemble $\underbrace{ A \times A \times \cdots  \times A }_{k-\text{occurences}} $ des $k$-uplets d’éléments d’un ensemble $A$. \\

$$ A^k=\{(a_1;a_2;\cdots;a_k)/ a_i \in A \; \text{pour tout} \; i\in \llbracket 1;n \rrbracket \} $$
\end{defn}

\subsection{Cardinal d'un ensemble}
\begin{defn}
La cardinalité est une notion de taille pour les ensembles, il s'agit du nombre d'éléments de l'ensemble. On note $\#A$ le \textbf{cardinal} de $A$.
\end{defn}

\begin{rem}
En particulier, le cardinal de l'ensemble vide est zéro. $(\#\emptyset =0)$
\end{rem}

\begin{prop}
Soit $A\subset B$ alors $\#A \ie \#B$.
\end{prop}

\subsection{Opérations sur les ensembles et cardinalité}
\begin{prop}
Le cardinal de la \textbf{réunion} d'ensemble: $ \# A\cup B \ie \#A +\#B$ .
\par Le cardinal de la \textbf{réunion disjointe} d'ensemble: $ \# A\uplus B = \#A +\#B$
\par Le cardinal de l'\textbf{intersection} d'ensemble: $ \# A\cap B \ie max(\#A,\#B)$.
\par Le cardinal du \textbf{produit cartésien} d'ensemble: $\# (A\times B)= \#A \times \#B$.
\end{prop}

\subsection{Nombre de $k$-uplets d'un ensemble à $n$ éléments}
\begin{prop}
Le nombre  de $k$-uplets d'un ensemble à $n$ éléments est égal à $n^k$ 
\end{prop}

\subsection{Nombre de $k$-uplets d'éléments distincts d'un ensemble à $n$ éléments}
\begin{prop}
Le nombre  de $k$-uplets d'un ensemble à $n$ éléments est égal à $n(n-1)(n-2)\cdots(n-k)$. 
\end{prop}

\subsection{Nombre de sous ensemble d’un ensemble à $n$ éléments}
Soit $A$ un ensemble fini à $n$ éléments que l'on note $a_1,a_2,\cdots,a_n$. Soit $B$ un sous ensemble de $A$, à $B$, on associe un $n$-uplet de la manière suivante.
\begin{itemize}
\item[$\bullet$]Dans la $i$-ième coordonnée, mettre 1 si $a_i \in B$ et 0 sinon.
\end{itemize}
Ainsi n'importe quel sous-ensemble est identifié avec un \textbf{unique} $n$-uplet de l'ensemble $\{0;1\}$, ainsi d'apr\`es la propriété sur la cardinalité on a $\# \{0;1\}^n =2^n$ sous ensemble d'un ensemble à $n$.  
\section{Permutations}

Le Père Noël a offert à ma petite cousine un jeu de cubes
où sont inscrits les lettres de l'alphabet.\\

Très pédagogue, je lui donne d'abord les trois cubes A, B et C.
Combien de " mots " de 3 lettres peut-elle alors former~? Et si
je lui en donne quatre~? vingt-six~? trente-deux~?\\


Moralité~: avec un ensemble \textit{(non ordonné)} $\{a,b,c\}$ de
trois éléments, je peux former $3\times 2\times 1$ listes
\textit{(ordonnées)}, comme par exemple $(a,b,c)$, $(b,a,c)$,
$(c,b,a)$, etc.

Je peux donc \textit{permuter} 3 cubes de $3\times 2\times 1$
manières différentes.

\begin{defn}[factorielle]
  Soit $n$ un entier naturel non nul. Le nombre $n\times(n-1)\times\cdots\times2\times1$ est appelé \textbf{factorielle n}. On
note
$$n!=n\times(n-1)\times\cdots\times2\times1$$
Conventionnellement, $0!=1$.
\end{defn}



D'après notre étude, nous pouvons donc dire maintenant que~:


\begin{thm}[nombre de permutations]
  Le nombre de \textbf{permutations} d'un ensemble contenant $n$ éléments est $n!$.
\end{thm}


\begin{exemple}[]
  \begin{enumerate}
  \item il y a trente-deux chaises dans une salle.
\textit{Dénombrer} toutes les manières possibles pour les 30 élèves
de $TS_4$ d'occuper les chaises de la salle~;
\item Combien y a-t-il d'annagrammes du mot ZO\'E~? du mot ANA~?
  \end{enumerate}
\end{exemple}


\section{Combinaisons}

Il n'y a pas assez de gagnants au loto. Les règles en sont donc
modifiées. Il s'agit maintenant de trouver 3 numéros parmi 5.
Combien y a-t-il de grilles (combinaisons) possibles~?

Par exemple, on pourrait dire que j'ai cinq manières de choisir le
premier numéro, quatre choix pour le deuxième et trois choix pour
le troisième, donc il y a $5\times 4\times 3$ grilles différentes.\\

Posons une petite définition pour clarifier les débats. Donnons en
fait un nom à une grille du loto, c'est à dire à un sous-ensemble
(une partie) contenant $k$ éléments d'un plus grand ensemble
contenant $n$ éléments.

\begin{defn}[combinaison]
  Soit $n$ et $k$ deux entiers naturels et $E$ un ensemble contenant $n$ éléments. Un sous-ensemble de $E$ contenant $k$
éléments est appelé une \textbf{combinaison} de $k$ éléments de $E$ ou encore une $k$\textbf{-combinaison} d'éléments de
\unboldmath${E}$.
\end{defn}


Or ce qui nous intéresse, c'est le nombre de ces combinaisons,
donc introduisons une notation~:
\subsection{Nombre de sous ensemble à $k$ éléments parmis un ensemble à $n$ éléments}
\begin{defn}[nombre de combinaisons]
   Le nombre de sous ensemble à $k$ éléments (\textit{ appelé aussi parfois $k$-combinaisons}) parmi un ensemble à $n$ éléments est $\co{k}{n}$. On le prononcera le nombre de combinaison de $k$ parmis $n$
  \vspace{3mm}
\end{defn}

Revenons à notre mini-loto. Considérons une grille quelconque
(c'est à dire une 3-combinaison de l'ensemble des 5 numéros)~: par
exemple $\{2,4,5\}$. Nous  avons vu dans le paragraphe précédent
qu'il y a $3!$ facons d'ordonner ces nombres. Finalement, il y a
$\co{3}{5} \times 3!$ suites de 3 nombres ordonnées. Or nous en
avons comptées $5\times 4\times 3$ tout à l'heure. Nous en
déduisons finalement que
$$\co{3}{5}  =\dfrac{5\times 4\times 3}{3!}$$

Il est alors aisé de généraliser la formule suivante~:


\begin{prop}[]
  $$\co{k}{n}=\dfrac{n(n-1)(n-2)(n-3)\cdots(n-k)}{k!}$$
\end{prop}



Nous pouvons formuler cette propriété plus synthétiquement. En
effet~:




\begin{thm}[]
   $$\co{k}{n}=\dfrac{n!}{k!(n-k)!}$$
\end{thm}

\begin{demo}
\begin{eqnarray*}
  {\co{k}{p}}&=&\dfrac{n(n-1)(n-2)\cdots(n-k+1)}{k!}\times
\dfrac{(n-k)(n-k-1)(n-k-2)\times\cdots\times 2\times
1}{(n-k)(n-k-1)(n-k-2)\times\cdots\times 2\times
1}\\
&=&\dfrac{n!}{k!(n-k)!}
\end{eqnarray*}
\end{demo}


\begin{prop}
$\sum^n_{k=1} \co{k}{n} =2^n$
\end{prop}

\begin{demo}

On sait que le nombre total de sous-ensemble d'un ensemble $n$ éléments est $2^n$, de plus le nombre de sous-ensemble à $k$ éléments parmi un ensemble à $n$ éléments est égal à $\co{k}{n}$. \\ 

Mais le nombre total de sous-ensemble d'un ensemble $n$ éléments ($2^n$) est la somme de tous les sous-ensembles à $1$ élément parmi un ensemble à $n$ éléments $(\co{1}{n})$, à $2$ éléments parmi un ensemble à $n$ éléments  $(\co{2}{n})$, ... , à $n$ éléments parmi un ensemble à $n$ éléments  $ (\co{n}{n})$. 

Autrement dit, 

$$ 2^n= \co{1}{n} + \co{2}{n} + \cdots +  \co{n}{n} $$

ou encore 
$$\sum^n_{k=1}  \co{k}{n} = 2^n $$
\end{demo}
 \newpage
\section{Triangle de Pascal - Binôme de Newton}

\subsection{Raisonnement calculatoire}

À l'aide des formules précédentes, établissez que~:


\begin{prop}[Symétrie]
$$\co{k}{n}=\co{{n-k}}{n}$$
\end{prop}


\'Etablissez, toujours par le calcul, la relation suivante, dite \textbf{Relation de Pascal} même si les mathématiciens chinois l'avaient
mise en évidence avant lui~:


\begin{prop}[Relation de \textsc{Pascal}]
   $$\co{k}{n}=\co{{k}}{{n-1}}+\co{{k-1}}{{n-1}}$$
\end{prop}

\begin{demo}
\begin{align*}
\co{{k}}{{n-1}}+\co{{k-1}}{{n-1}} &=\dfrac{(n-1)!}{(n-1-k)!k!}+\dfrac{(n-1)!}{(n-k)!k-1!}\\
& = \dfrac{(n-k)(n-1)!}{(n-k)!k!}+\dfrac{k(n-1)!}{(n-k)!k!} \\
& = \dfrac{(n-k)(n-1)! + k(n-1)!}{(n-k)!k!} \\
& = \dfrac{(n-k + k)(n-1)!}{(n-k)!k!} \\
& =\dfrac{n(n-1)!}{(n-k)!k!} \\
& =\dfrac{n!}{(n-k)!k!}\\
 & =:\co{k}{n} \; \text{par définition}
\end{align*}
d'où 
$$\co{k}{n}=\co{{k}}{{n-1}}+\co{{k-1}}{{n-1}}$$

\end{demo}



Cette relation permet de calculer les coefficients binomiaux de proche en proche, ce qui s'avère fort utile car la fonction factorielle
croît extrêmement rapidement et dépasse vite les capacités d'une calculatrice de poche.

Ainsi, la procédure suivante permet de calculer les coefficients binomiaux assez rapidement avec

% \Prog{XCAS}

%\begin{lstlisting}
%combi(p,n):={
 %si p=0 alors 1; 
 %sinon si p>n alors 0;
    %   sinon combi(p,n-1)+combi(p-1,n-1);
      % fsi
% fsi
%}:;
%\end{lstlisting}


%Mais que viennent faire les deux conditions : \texttt{si p=0} et \texttt{si p>n}~?
\newpage
\subsection{Raisonnement ensembliste}

Si je forme un groupe de 3 élèves dans la classe, j'obtiens du même coup un groupe de 26 élèves puisque vous êtes 29.

Si l'on compte le nombre de parties A ayant $p$ éléments dans un ensemble E, il revient au même de compter le nombre de parties
complémentaires de A. En quoi ce raisonnement nous permet-il d'obtenir une des formules précédentes~?




Considérons la classe de TS4 et ses 30 élèves. Je veux dénombrer les groupes de trois élèves. On peut distinguer deux types de
groupes~:

\begin{itemize}
\item ceux qui ne vous contiennent pas : il faut donc former des groupes de 3 avec les 22 élèves restant : ça fait combien~?
\item ceux qui vous contiennent : on vous ajoute aux groupes de 2 formés avec les 22 élèves restant. Ça fait combien~?
\end{itemize}



Tenez   un  raisonnement   similaire   pour  établir   la  relation   de
\textsc{Pascal}.



\subsection{Formule du binôme de \textsc{Newton}}

Sir Isaac  n'a pas  fait que dormir  sous un  pommier. Il a  par exemple
établi que~:


\begin{thm}[Formule du binôme de \textsc{Newton}]
   Soit $a$ et $b$ des nombres réels ou complexes et $n$ un entier naturel non nul, alors
$$(a+b)^n=\sum_{k=0}^n\co{k}{n}a^kb^{n-k}=a^n+\co{1}{n}a^{n-1}b+\cdots+\co{{n-1}}{n}ab^{n-1}+b^n$$
\end{thm}


On peut prouver cette formule par récurrence en remarquant que $(a+b)^{k+1}=a(a+b)^k+b(a+b)^k$ et en
utilisant la relation de Pascal au bon moment.\\ 

Cette formule nous permet donc d'obtenir de nouveaux " produits remarquables ", à conditions de connaître les coefficients binomiaux.

Testez la formule aux rangs 2, 3, 4, 5. Disposez vos résultats dans un tableau en n'écrivant que les coefficients et conjecturer le
triangle de \textsc{Pascal}...

\begin{center}
\begin{tabular}{>{$}l<{$}|*{7}{c}}
\multicolumn{1}{l}{$k$} &&&&&&&\\\cline{1-1} 
0 &1&&&&&&\\
1 &1&1&&&&&\\
2 &1&2&1&&&&\\
3 &1&3&3&1&&&\\
4 &1&4&6&4&1&&\\
5 &1&5&10&10&5&1&\\
6 &1&6&15&20&15&6&1\\\hline
\multicolumn{1}{l}{} &0&1&2&3&4&5&6\\\cline{2-8}
\multicolumn{1}{l}{} &\multicolumn{7}{c}{$n$}
\end{tabular}
\end{center}
\begin{comment}

\begin{center}
\includegraphics[height=8cm]{triangle}
\end{center}




\modeexercice % entre en mode exercice
\section{EXERCICES}


\begin{listeexercices}

  \begin{exercice}[Coefficients binômiaux]
    \begin{enumerate}
    \item Donnez une expression simple de $\binom{n}{0}$, de $\binom{n}{1}$, de $\binom{n}{n}$, de $\binom{n}{n-1}$. Utilisez des
simplifications pour calculer à la main $\binom{20}{3}$.

\item On donne $\binom{13}{5}=1287$ et $\binom{13}{6}=1716$. Calculez alors à la main $\binom{13}{8}$ et $\binom{14}{9}$
    \end{enumerate}
    
  \end{exercice}


  \begin{exercice}[Une vieille connaissance]
    
Donnez une nouvelle démonstration très rapide de la formule bien connue : $(1+x)^n>1+nx$

  \end{exercice}


  \begin{exercice}[Nombre de parties]
    
Combien y a-t-il de parties d'un ensemble ayant $n$ éléments~?

  \end{exercice}


  \begin{exercice}[Bac avec ROC]
    

\begin{enumerate} \item \textbf{Démonstration de cours.} Démontrer que, pour tous entiers naturels
$n$ et $k$ tels que $1 \leqslant  k < n$, on a :

\[\binom{n-1}{k - 1} + \binom{n - 1}{k} = \binom{n}{k}.\]

\item En déduire que pour tous entiers naturels $n$ et $k$ tels que $2 \leqslant k <
n - 1$, on a :
\[\binom{n-2}{k-2} + 2 \binom{n-2}{k-1} + \binom{n-2}{k} =
\binom{n}{k}.\]

\item On considère deux entiers naturels $n$ et $k$ tels que $2 \leqslant k <
n - 1$. On dispose d'une urne contenant $n$ boules indiscernables au toucher. Deux des boules sont rouges, les autres sont blanches.

On tire au hasard et simultanément $k$ boules de l'urne. On appelle A l'évènement « au moins une boule rouge a été tirée ».

\begin{enumerate} \item Exprimer en fonction de $n$ et de $k$ la probabilité
de l'évènement $\overline{\text{A}}$, contraire de A.

En déduire la probabilité de A.

 \item Exprimer d'une autre manière la probabilité de l'évènement A et
montrer, à l'aide la formule obtenue à la question \textbf{2.}, que l'on retrouve le même résultat.

\end{enumerate}

\end{enumerate}
  \end{exercice}




  \begin{exercice}
    Calulez $\ds\int_{\efr{\pi}{6}}^{\efr{\pi}{3}}\sin^5t \mathrm{d}t$

  \end{exercice}


  \begin{exercice}[exercice grec]
    \begin{center}\includegraphics[width=0.7\linewidth]{ouzo}\end{center}

 Périclès est goutteur d'olives dans une usine grecque. Un
matin, il goutte cent olives au hasard et les replace dans le réservoir.

L'après-midi, l'ouzo de l'apéritif lui a fait perdre la mémoire. Il goutte à nouveau cent olives dans le même réservoir. Douze d'entre
elles avaient déjà été mâchées.

On note $A$ l'événement \og il y a douze olives mâchées parmi les cent choisies \fg{} et $B_n$ l'événement \og il y a $n$ olives dans le
réservoir \fg{}.

On considère la fonction $f$ définie pour les entiers supérieurs à cent  par
$f(n)=\bbp_{B_n}(A)$ et la suite $(u_n)$ définie pour les entiers
supérieurs à 100 par
$u_n=f(n+1)/f(n)$.
\been \item Comparez $u_n$ à 1.
\jok{$u_n=\fr{\binom{100}{12}\binom{n-100}{88}\binom{n+1}{100}}{\binom{100}{12}\binom{n-99}{88}\binom{n}{100}}$}
 \item Montrez que la fonction $f$ atteint un maximum sur $[\![100,+\infty[\![$. \item On appelle
\textsl{maximum de vraissemblance} $m$ la valeur de $n$ correspondant à ce maximum. Déterminez $m$. \een
\begin{flushright}\rotatebox{180}{Réponse~: $m=833$}\end{flushright}
  \end{exercice}


\end{listeexercices}


\modenormal

%%% Local Variables: 
%%% mode: latex
%%% TeX-master: "/home/moi/Lycee/TS/2009_10/PolyTs/PolyTs.tex"
%%% End: 
\end{comment}
\end{document}
